\[ %%%%%%%%%%%%%%%%%%%%%%%%%%%%%%%%%%%%%%%%%%%%%%%%%%%%%%%%%%%%%%%%%%%%%%%%%%%%%%%% \subsection{Generalization Across Physical Systems} \label{sec:generalization} %%%%%%%%%%%%%%%%%%%%%%%%%%%%%%%%%%%%%%%%%%%%%%%%%%%%%%%%%%%%%%%%%%%%%%%%%%%%%%%% A fundamental objective of PHYSGEN is to develop a general framework capable of representing and predicting a wide range of nonlinear physical systems. Unlike domain-specific simulation models that are designed for a particular equation set or application area, PHYSGEN aims to learn general principles of physical evolution through adaptive graph representations and physics-guided dynamics. %%%%%%%%%%%%%%%%%%%%%%%%%%%%%%%%%%%%%%%%%%%%%%%%%%%%%%%%%%%%%%%%%%%%%%%%%%%%%%%% \subsubsection{Unified Representation of Physical Systems} %%%%%%%%%%%%%%%%%%%%%%%%%%%%%%%%%%%%%%%%%%%%%%%%%%%%%%%%%%%%%%%%%%%%%%%%%%%%%%%% Although physical systems differ significantly in their underlying mechanisms, many of them share a common mathematical structure: \begin{equation} X_{t+1} = \mathcal{F}(X_t,\Theta). \end{equation} PHYSGEN assumes that the essential information required for prediction can be represented through: \begin{equation} G_t=(V_t,E_t,X_t,P_t), \end{equation} where: \begin{itemize} \item nodes represent physical entities or spatial locations; \item edges represent interactions or dependencies; \item node features represent physical states; \item parameters represent domain-specific properties. \end{itemize} This representation allows heterogeneous physical systems to be mapped into a common computational framework. %%%%%%%%%%%%%%%%%%%%%%%%%%%%%%%%%%%%%%%%%%%%%%%%%%%%%%%%%%%%%%%%%%%%%%%%%%%%%%%% \subsubsection{Domain-Invariant Physical Operators} %%%%%%%%%%%%%%%%%%%%%%%%%%%%%%%%%%%%%%%%%%%%%%%%%%%%%%%%%%%%%%%%%%%%%%%%%%%%%%%% The central hypothesis of PHYSGEN is that physical evolution can be decomposed into two components: \begin{equation} \Phi = \Phi_{universal} + \Phi_{domain}. \end{equation} The universal component captures general dynamical principles: \begin{itemize} \item locality of interactions; \item conservation behavior; \item temporal evolution; \item stability conditions. \end{itemize} The domain-specific component captures: \begin{itemize} \item material properties; \item boundary conditions; \item system parameters; \item application-specific constraints. \end{itemize} Therefore, PHYSGEN can adapt between physical domains while preserving learned dynamical structures. %%%%%%%%%%%%%%%%%%%%%%%%%%%%%%%%%%%%%%%%%%%%%%%%%%%%%%%%%%%%%%%%%%%%%%%%%%%%%%%% \subsubsection{Cross-Domain Transfer Formulation} %%%%%%%%%%%%%%%%%%%%%%%%%%%%%%%%%%%%%%%%%%%%%%%%%%%%%%%%%%%%%%%%%%%%%%%%%%%%%%%% Let a physical domain be defined as: \begin{equation} \mathcal{D}_i = \{ G_i,P_i,\mathcal{C}_i \}. \end{equation} The objective is to learn parameters: \begin{equation} \theta = \{\theta_{shared},\theta_i\}, \end{equation} where: \begin{itemize} \item $\theta_{shared}$ represents transferable physical knowledge; \item $\theta_i$ represents domain-specific adaptation. \end{itemize} The prediction process becomes: \begin{equation} \hat{X}_{t+1} = \Phi_{\theta_{shared},\theta_i} ( G_t,P_i,\mathcal{C}_i ). \end{equation} This enables transfer learning between related physical systems. %%%%%%%%%%%%%%%%%%%%%%%%%%%%%%%%%%%%%%%%%%%%%%%%%%%%%%%%%%%%%%%%%%%%%%%%%%%%%%%% \subsubsection{Examples of Applicable Physical Domains} %%%%%%%%%%%%%%%%%%%%%%%%%%%%%%%%%%%%%%%%%%%%%%%%%%%%%%%%%%%%%%%%%%%%%%%%%%%%%%%% The general formulation of PHYSGEN allows application to multiple classes of dynamical systems. \paragraph{Fluid Dynamics.} For computational fluid dynamics, nodes represent spatial discretization points and edges describe local flow interactions: \begin{equation} G_{fluid} = (V_{space},E_{flow},X_{velocity,pressure}). \end{equation} \paragraph{Material and Molecular Systems.} For molecular or materials simulations: \begin{equation} G_{material} = (V_{atoms},E_{bonds},X_{position,energy}). \end{equation} \paragraph{Biological Systems.} For biological digital twins: \begin{equation} G_{bio} = (V_{components},E_{interaction},X_{state}). \end{equation} \paragraph{Climate and Earth Systems.} Large-scale environmental systems can be represented as: \begin{equation} G_{climate} = (V_{regions},E_{transport},X_{fields}). \end{equation} %%%%%%%%%%%%%%%%%%%%%%%%%%%%%%%%%%%%%%%%%%%%%%%%%%%%%%%%%%%%%%%%%%%%%%%%%%%%%%%% \subsubsection{Toward Physics Foundation Models} %%%%%%%%%%%%%%%%%%%%%%%%%%%%%%%%%%%%%%%%%%%%%%%%%%%%%%%%%%%%%%%%%%%%%%%%%%%%%%%% The ability to transfer learned physical representations across domains suggests a path toward physics foundation models. A physics foundation model can be formulated as learning a universal operator: \begin{equation} \Phi_{\Theta}^{*} : \mathcal{G} \rightarrow \mathcal{G}, \end{equation} where $\mathcal{G}$ represents the space of physical graphs. The goal is not to memorize individual simulations, but to learn reusable principles of physical evolution. PHYSGEN provides several components required for such a framework: \begin{itemize} \item graph-based physical representation; \item physics-guided latent dynamics; \item stability-aware evolution; \item multi-level physical constraints. \end{itemize} %%%%%%%%%%%%%%%%%%%%%%%%%%%%%%%%%%%%%%%%%%%%%%%%%%%%%%%%%%%%%%%%%%%%%%%%%%%%%%%% \subsubsection{Generalization Objective} %%%%%%%%%%%%%%%%%%%%%%%%%%%%%%%%%%%%%%%%%%%%%%%%%%%%%%%%%%%%%%%%%%%%%%%%%%%%%%%% The generalization objective can be formulated as minimizing expected error across multiple physical domains: \begin{equation} \theta^* = \arg\min_\theta \sum_i \mathbb{E}_{\mathcal{D}_i} [ \mathcal{L}_{PHYSGEN} ( \mathcal{D}_i ) ]. \end{equation} This formulation encourages the model to discover common physical structures rather than overfitting to a single dataset. %%%%%%%%%%%%%%%%%%%%%%%%%%%%%%%%%%%%%%%%%%%%%%%%%%%%%%%%%%%%%%%%%%%%%%%%%%%%%%%% \subsubsection{Summary} %%%%%%%%%%%%%%%%%%%%%%%%%%%%%%%%%%%%%%%%%%%%%%%%%%%%%%%%%%%%%%%%%%%%%%%%%%%%%%%% The generalization formulation positions PHYSGEN as a universal framework for learning nonlinear physical dynamics across multiple scientific domains. By representing different systems as evolving physical graphs and integrating domain-independent physical principles, PHYSGEN provides a foundation for scalable scientific machine learning and future physics foundation models. \] 